\addcontentsline{toc}{subsubsection}{Solution 4. The Dyadic Lock}

\begin{dlsummary}[title={\bfseries Solution 4. The Dyadic Lock},breakable]

\begin{proof}
The second inequality in the problem implies
\[
f(xy)\leq f(x)f(y). \tag{1}
\]
Also, $f$ is strictly increasing: if $u>v>0$, then
$f(u)=f(v+(u-v))\geq f(v)+f(u-v)>f(v)$.
Taking $x=y=1$ in (1) gives $f(1)\geq1$. Superadditivity then gives
$f(2)\geq2f(1)\geq2$. Since $f$ is increasing,
$f(f(2))\geq f(2)$, and therefore
\[
4\geq f(2)f(f(2))\geq f(2)^2\geq4.
\]
It follows that $f(2)=2$ and, from $f(2)\geq2f(1)\geq2$, also
$f(1)=1$. Superadditivity and (1), applied repeatedly, now give
\[
f(2^n)=2^n\qquad(n=0,1,2,\ldots). \tag{2}
\]

We next obtain a common upper bound. Fix $t>0$, and for all sufficiently
large $n$ let
\[
m_n=\left\lfloor\frac{2^n}{t}\right\rfloor.
\]
By monotonicity, superadditivity and (2),
\[
2^n=f(2^n)\geq f(m_nt)\geq m_nf(t).
\]
Since $2^n/m_n\to t$, this proves
\[
f(t)\leq t\qquad(t>0). \tag{3}
\]

For every positive integer $k$, superadditivity and (3) give
\[
k=kf(1)\leq f(k)\leq k,
\]
so $f(k)=k$. If $r=p/q$ is a positive rational, where $p,q$ are
positive integers, then (1) yields
\[
p=f(p)=f(qr)\leq f(q)f(r)=qf(r).
\]
Hence $f(r)\geq r$, and (3) gives $f(r)=r$. For any positive real $x$,
choose positive rationals $r<x<s$ arbitrarily close to $x$. Monotonicity
then gives $r=f(r)\leq f(x)\leq f(s)=s$. Letting $r,s$ approach $x$
proves $f(x)=x$.

The identity function satisfies the original inequalities and numerical
condition with equality, so it is the unique solution.
\end{proof}

\end{dlsummary}